\footnotesize
\arrayrulecolor{black!22}
\begin{longtable}{L{0.04\linewidth} L{0.105\linewidth} L{0.16\linewidth} L{0.13\linewidth} L{0.125\linewidth} L{0.165\linewidth} L{0.075\linewidth} L{0.135\linewidth}}
\hline
\textbf{ID} & \textbf{Topic} & \textbf{Source evidence} & \textbf{External check} & \textbf{Plan disposition} & \textbf{CNNA action} & \textbf{Phase/\allowbreak{}Object} & \textbf{Risk /\allowbreak{} uncertainty}\\
\hline
\endfirsthead
\hline
\textbf{ID} & \textbf{Topic} & \textbf{Source evidence} & \textbf{External check} & \textbf{Plan disposition} & \textbf{CNNA action} & \textbf{Phase/\allowbreak{}Object} & \textbf{Risk /\allowbreak{} uncertainty}\\
\hline
\endhead
TOMS-01 & Prime generation vs irreducible/\allowbreak{}prime distinction & \#8-\#10 and \#48: Schur levels, Eisenstein/\allowbreak{}Gauss norm heuristics, TomS correction that irreducible need not be prime & Algebraically standard: UFD/\allowbreak{}PID/\allowbreak{}class-number gates are required before irreducible=prime can be claimed & Already broadly covered by NumberTheory phases but wording must stay strict. & Keep Heegner/\allowbreak{}PID/\allowbreak{}UFD indicator after class number;\newline never claim prime generation from norm irreducibility without class-number/\allowbreak{}UFD proof. & 6.2, 6.3 /\allowbreak{} N09-N13 & The v3 script observations remain empirical/\allowbreak{}heuristic until operational recurrence, norm form, and class group gates are proved.\\
\hline
\addlinespace[0.24em]
TOMS-02 & Quadratic orders and lattice types & \#9-\#11: imaginary-quadratic integer rings have two order-shapes by d mod 4;\newline elliptic functions exist on arbitrary non-collinear complex lattices & Established CM theory: quadratic orders/\allowbreak{}lattices are restrictive subfamily of all elliptic lattices & Covered by D\_fund/\allowbreak{}order/\allowbreak{}lattice/\allowbreak{}tau phases;\newline add caution to source map. & Do not say Schur levels generate all elliptic lattices;\newline derive D/\allowbreak{}order/\allowbreak{}basis/\allowbreak{}tau first, then compare Gauss/\allowbreak{}Eisenstein special cases. & 5.1-8.3 /\allowbreak{} N00-N07, L00-L07 & Current code has special Gaussian/\allowbreak{}Eisenstein bridges but no general D\_op->order pipeline.\\
\hline
\addlinespace[0.24em]
TOMS-03 & SL2Z/\allowbreak{}Moebius, modular functions, j & \#13/\allowbreak{}\#15/\allowbreak{}\#22: SL2Z matrices act by z -> (az+b)/\allowbreak{}(cz+d), S and T generate transformations;\newline j as invariant modular function & Milne/\allowbreak{}standard modular-form references support j as generator of modular-function field;\newline thread contains minor typo-like fundamental-domain bound, so do not copy formulas unchecked & Already present but should remain after CM/\allowbreak{}tau, not before. & Keep modular group action, fundamental-domain representative, q-expansion and Hilbert class polynomial as post-lattice/\allowbreak{}modular phases. & 9, 10.3 /\allowbreak{} M00-M06 & Need authoritative source and mathlib import check at implementation;\newline do not treat forum formulas as formal definitions.\\
\hline
\addlinespace[0.24em]
TOMS-04 & j algebraic at CM points & \#30: j(tau) algebraic on imaginary-quadratic tau;\newline CM connection & Milne modular-function notes and CM theory support algebraicity at imaginary-quadratic points & Covered by CM/\allowbreak{}lattice/\allowbreak{}j phases but dependency should be explicit. & Require tau upper-half-plane and CM/\allowbreak{}order provenance before j-values/\allowbreak{}class-polynomial work. & 8-10.3 /\allowbreak{} L04-L05, M05-M06 & Exact field/\allowbreak{}class-polynomial statements require class group/\allowbreak{}ring class field layer;\newline current plan only has placeholders.\\
\hline
\addlinespace[0.24em]
TOMS-05 & L-functions are not Riemann-zeta-first & \#49: TomS says zeta/\allowbreak{}L-functions yes, not specifically Riemann zeta;\newline modular forms are the bridge & Standard theory: many L-functions attach to characters/\allowbreak{}modular forms;\newline Riemann zeta is special with trivial character & Plan should not introduce Riemann-zeta as generator. & Keep chi\_D/\allowbreak{}splitting before any L-function comparison;\newline any L-function object is derived/\allowbreak{}comparison after character/\allowbreak{}modular data. & 6.4, 10.5 /\allowbreak{} N15, spectral-zeta side object & Avoid conflating spectral zeta, Dirichlet L-functions, and Riemann zeta.\\
\hline
\addlinespace[0.24em]
TOMS-06 & Finite simple groups taxonomy & \#43-\#56: cyclic, alternating, Lie type, sporadic groups;\newline groups as generalized prime-like objects & Warwick WMI: CFSG categories are cyclic prime order, alternating, Lie type, sporadic & Gap in v20: Moonshine section lacked explicit finite-simple-group taxonomy as comparison-only. & Add MOON3/\allowbreak{}MOON4 and phase 12.4 as external comparison context only. & 12.4 /\allowbreak{} MOON3-MOON4 & CFSG is extremely deep;\newline never import it into A proof path as an unchecked theorem.\\
\hline
\addlinespace[0.24em]
TOMS-07 & Monster is not final object;\newline common indicator matters & \#64-\#71: Monster/\allowbreak{}Moonshine;\newline \#67/\allowbreak{}\#70 says Monster is only tip of iceberg, seek connecting structure and common indicator & Borcherds 1992 proves Conway-Norton moonshine conjectures using vertex algebras, Thompson series/\allowbreak{}Hauptmoduls, Monster simple group & v20 had Monster/\allowbreak{}McKay-Thompson placeholders;\newline needs common-indicator/\allowbreak{}Hauptmodul bridge-test record. & Add MOON8 and phase 12.8;\newline make Monster comparison subordinate to derived q-character/\allowbreak{}Hauptmodul tests. & 12.8 /\allowbreak{} MOON8 & Do not overfit CNNA to Monster;\newline use it as late fingerprint test only.\\
\hline
\addlinespace[0.24em]
TOMS-08 & Leech lattice/\allowbreak{}rootlessness/\allowbreak{}theta shell data & \#84/\allowbreak{}\#93: Leech lattice has no roots;\newline shell/\allowbreak{}theta coefficients expose missing first shell and 196560 minimal vectors & Borcherds 1985 proves Leech lattice existence/\allowbreak{}uniqueness and records no-root/\allowbreak{}Niemeier context & Gap in v20: no Leech-specific comparison object. & Add MOON5/\allowbreak{}MOON6 and phases 12.5-12.6. & 12.5-12.6 /\allowbreak{} MOON5-MOON6 & Thread details around root systems were corrected by Jakito/\allowbreak{}TomS;\newline use corrected Erzeugendensystem not basis wording.\\
\hline
\addlinespace[0.24em]
TOMS-09 & VOA/\allowbreak{}chiral CFT/\allowbreak{}string-theory context & \#64: Moonshine connection to bosonic string compactified on Leech lattice and 2D chiral CFT;\newline \#82/\allowbreak{}\#85 caution about physics-from-math & Borcherds 1992 keywords include vertex algebras, monster simple group, Thompson series, Hauptmoduls & Needs strict comparison/\allowbreak{}consumer boundary. & Add MOON7 and keep string theory in side-comparison phase 18.2;\newline no VOA/\allowbreak{}string axiom import into Pillar A. & 12.7, 18.2 /\allowbreak{} MOON7 & VOA formalization is much deeper than current A path;\newline only finite graded truncations may be computable comparison artifacts.\\
\hline
\addlinespace[0.24em]
TOMS-10 & Elliptic functions are related but not the final bridge & \#82/\allowbreak{}\#89: TomS did not use elliptic functions as over-arching link;\newline Jakito notes G\_\{2k\}, Delta, j arise from elliptic-function lattice theory & Standard modular/\allowbreak{}elliptic-function pedagogy supports G4/\allowbreak{}G6 -> g2/\allowbreak{}g3 -> Delta -> j & Already covered by modular forms and Weierstrass invariant phases. & Keep elliptic-function material as necessary j-construction route, not as claimed universal CNNA endpoint. & 9.1-9.2 /\allowbreak{} M02-M05 & At implementation, distinguish analytic exact layer from finite q/\allowbreak{}truncation mirror.\\
\hline
\addlinespace[0.24em]
TOMS-11 & Fractal fundamental-domain imagery is low priority & \#22/\allowbreak{}\#24: fractal image language for iterated fundamental domains is explicitly uncertain/\allowbreak{}irrelevant in thread & No primary formal target needed for current CNNA plan & No action in near path. & Do not create a CNNA object for fractality;\newline if ever needed, classify as visualization or tessellation comparison after modular action exists. & 10.3 only if needed & Avoid spending proof effort on metaphorical/\allowbreak{}visual claims.\\
\hline
\addlinespace[0.24em]
TOMS-12 & Physics-from-math caution & \#37/\allowbreak{}\#85: TomS objects that physics is not obtained by pure mathematics alone;\newline thread is not a CNNA ontological source & Methodological warning aligns with comparison-only and handoff discipline & Supports existing access policy. & Keep established physics/\allowbreak{}math theories as references and tests, never as ontic/\allowbreak{}generative inputs. & access\_policy;\newline 18.* comparisons & Empirical/\allowbreak{}physical claims remain speculative unless reconstructed from CNNA and compared to experiment.\\
\hline
\addlinespace[0.24em]
TOMS-13 & Heuristic b-ary tree construction of Eisenstein/\allowbreak{}Gauss lattice profiles & post \#97: b-ary address construction;\newline b>=2 heuristic;\newline Eisenstein angle 120 degrees and closure of 0-path at level 3;\newline Gauss angle 90 degrees and closure at level 4;\newline hierarchical overlays versus spiral angle dependence. & The thread entry is heuristic and self-authored, not a theorem source. Established CM/\allowbreak{}j special targets should be verified from formal references during M05/\allowbreak{}M08;\newline forum data may define tests but not proofs. & Add green boundary lock 1.1.1.1.0.9 and future phases 8.5-8.7 for static/\allowbreak{}dynamic test matrices and closeout. & Use code-defined branching b. Test rows must cover angle, rational cosine, finite address collisions/\allowbreak{}closure, CM point target, discriminant target, and j target over selected L\_max values. Keep complex coordinates/\allowbreak{}tau/\allowbreak{}j analytic evaluation mirror-only under EXEC-2. & 1.1.1.1.0.9;\newline 8.5-8.7 /\allowbreak{} LEG28, L09, L10, M08, BR10 & The geometric embedding is not yet derived from the generator;\newline it may fail or remain comparison-only. Dynamic branching and general D require later source objects and cannot be assumed from b=2 drawings.\\
\hline
\addlinespace[0.24em]
TOMS14 & Post \#97 as early general test source, not special-case tunnel & Post \#97 motivates Eisenstein/\allowbreak{}Gauss angle/\allowbreak{}closure tests from b-ary tree drawings. & The source is heuristic and special-case oriented;\newline it does not define the code branching object or a complete general test suite. & Use as sentinel rows inside O07/\allowbreak{}L09/\allowbreak{}L10, not as the complete test matrix. & Add LEG29/\allowbreak{}O07 and GNG16;\newline update Phase 8.5/\allowbreak{}8.6 to require generic coverage before sentinel interpretation. & 1.1.1.1.0.10, 4.0.1, O07, L09/\allowbreak{}L10/\allowbreak{}M08/\allowbreak{}BR10 & If generic rows are too weak, sentinels may create false confidence. Coverage must be reviewed when real code starts.\\
\hline
\addlinespace[0.24em]
TOMS15 & Eisenstein/\allowbreak{}Gauss as noOuter-root sentinels, not the full Heegner story & Post \#97 gives b-ary address and angle/\allowbreak{}closure heuristics: 120 degrees with closure at L=3 for Eisenstein;\newline 90 degrees with closure at L=4 for Gauss;\newline the author also marks the discussion as heuristic in \#99. & Class-number-one/\allowbreak{}Heegner lists contain more than D=-3 and D=-4, so CNNA must not identify all Heegner targets with the closed/\allowbreak{}noOuter branch. Established lists are comparison targets, not CNNA derivations. & Add LEG30/\allowbreak{}O08/\allowbreak{}N18/\allowbreak{}BR11/\allowbreak{}BR12, GNG17, and phases 1.1.1.1.0.11, 4.0.2, 7.2.1. & Use regime-typed test rows: Eisenstein/\allowbreak{}Gauss are noOuter-root sentinel candidates;\newline other Heegner/\allowbreak{}CM targets are HasOuter/\allowbreak{}complement-projection or comparison rows unless noOuter provenance is derived. & 1.1.1.1.0.11;\newline 4.0.2;\newline 7.2.1 /\allowbreak{} LEG30, O08, N18, BR11, BR12 & This is a CNNA hypothesis and test classification, not yet a theorem. It may fail once concrete rows are computed;\newline failure should trigger backjump to regime/\allowbreak{}complement source layers, not a consumer patch.\\
\hline
\addlinespace[0.24em]
HORN-GA-01 & Horn 0709.2238v1 quaternion, Pauli and geometric-algebra synthesis & Uploaded PDF 0709.2238v1, sections on quaternions, Pauli matrices, geometric algebra and duality. & Horn presents quaternions, Pauli matrices and Hestenes geometric algebra as overlapping frameworks synthesized through duality;\newline the paper is didactic/\allowbreak{}conceptual and not a CNNA proof source. & Create comparison/\allowbreak{}reference gates GA0-GA3 under 1.4.9.* and keep them outside the derived-only generator chain. & Record target identities and ambiguity warnings: Hamilton unit products, Pauli/\allowbreak{}Cartan sign choices, complex-i versus quaternion-i, pseudoscalar duality, pure quaternion as bivector and biquaternion multivector bookkeeping. & 1.4.9;\newline 1.4.9.1;\newline 1.4.9.2 /\allowbreak{} GA0-GA3 & Sign conventions, orientation choices and semantics are underdetermined by the external reference;\newline a CNNA-derived object may fail the target or match only after a natural-origin proof supplies missing structure.\\
\hline
\addlinespace[0.24em]
HORN-GA-02 & Mathlib Quaternion and CliffordAlgebra seams as future formal target APIs & Web source review of current mathlib Quaternion and CliffordAlgebra documentation plus Lean geometric-algebra formalization literature. & Mathlib exposes Quaternion/\allowbreak{}QuaternionAlgebra and CliffordAlgebra APIs;\newline Wieser/\allowbreak{}Song note GA formalization through CliffordAlgebra and gaps around wedge/\allowbreak{}N-grading. & Record as possible future comparison implementation seams only;\newline no import is added before local Lean/\allowbreak{}mathlib 4.28.0 verification. & Add scratchpad expected-import warnings for 1.4.9.*, 13.2.1 and 13.6.1;\newline route implementation through minimal target profiles if mathlib imports are too heavy or noncomputable. & 1.4.9.1;\newline 1.4.9.2;\newline 13.2.1;\newline 13.6.1 /\allowbreak{} GA1-GA5 & Live docs may differ from the pinned project snapshot. Noncomputable or classical pressure must remain in mirror/\allowbreak{}comparison layers and never enter operational paths.\\
\hline
\addlinespace[0.24em]
HORN-GA-03 & Dirac, spacetime algebra and conformal GA as late target gates only & Horn discussion of quaternionic Dirac motivation, complex quaternions, outlook toward spacetime algebra and conformal geometric algebra. & The material is useful for target-gate design but cannot supply Lorentz structure, Dirac equations, spacetime dimension or conformal geometry as CNNA inputs. & Add 13.6.1 and 18.5 as late comparison phases with GA5-GA6. & Use after Pillar-B local-net/\allowbreak{}modular data and side-comparison outputs exist;\newline compare signatures and structural compatibility, not foundations. & 13.6.1;\newline 18.5 /\allowbreak{} GA5-GA6 & Potentially high heuristic pull toward established physics. Guardrail: no Dirac/\allowbreak{}STA/\allowbreak{}CGA theorem enters CNNA before derived predecessor objects exist.\\
\hline
\addlinespace[0.24em]
ROADMAP-01 & Roadmap source hierarchy and definition/\allowbreak{}notation discipline & Original roadmap sections Dokumenttyp/\allowbreak{}Quellenhierarchie, Definitionskatalog, Notationspflicht and Register state that the roadmap is a gate document, that every core object needs one canonical row, primary readable form, exact Lean binding, notation target module/\allowbreak{}scope and alias/\allowbreak{}register separation. & Internal roadmap evidence only;\newline it is not checked Lean. Its discipline aligns with current Single-YAML-master, object register and total declaration classification policies. & Add RDV0 and Phase 22.0;\newline tighten notation/\allowbreak{}register handling as a prerequisite for Explorer/\allowbreak{}source-intake closure. & Do not let glossary/\allowbreak{}register/\allowbreak{}notation rows create Fachobjekte. Route each non-meta form to one canonical object or explicit non-object bucket. & 22.0 /\allowbreak{} RDV0, CEX11, LEG37 & Risk: the original roadmap contains many planned/\allowbreak{}legacy names;\newline without total owner/\allowbreak{}bucket mapping it could inflate the object universe or duplicate aliases.\\
\hline
\addlinespace[0.24em]
ROADMAP-02 & Migration law: direct port, semantic rewrite, bright-special-case and late seed & Original roadmap migration-law sections distinguish portable kernels, modules requiring semantic re-reading, Bright-special-case regression material and late seeds/\allowbreak{}diagnostics. Boundary-first is explicitly not the new architecture center. & This is a migration audit rule, not a proof of any object. It matches the current natural-provenance patch rule and no phase-semantic public API naming policy. & Add RDV1 and Phase 22.1;\newline require migration disposition for legacy modules/\allowbreak{}objects before reuse or de-seeding. & Classify every legacy candidate as direct\_port, semantic\_rewrite, bright\_special\_case\_regression\_oracle or late\_seed\_or\_diagnostic;\newline keep old BoundaryMatrix path as regression/\allowbreak{}reference where appropriate. & 22.1 /\allowbreak{} RDV1 & Risk: direct porting can silently preserve obsolete Boundary-first semantics if disposition is not recorded before refactor.\\
\hline
\addlinespace[0.24em]
ROADMAP-03 & De-seeding, old-path regression, immediate readability and rollback & Original roadmap operative rules require SourceMap, removal of free cutoffs/\allowbreak{}states/\allowbreak{}dictionaries/\allowbreak{}observers, derived-only gate, notation binding, old-path regression oracle, pinned toolchain/\allowbreak{}build aggregators/\allowbreak{}importgraph audit and rollback artifacts. & These are governance conditions;\newline they do not close any Lean theorem by themselves. They complement object lifecycle and phase patch-link workflows. & Add RDV2, Phase 22.2, workflows and guardrail GNG46. & No Seed/\allowbreak{}Scaffold/\allowbreak{}Window module loses its suffix or becomes current API without source map, derived-only gate, notation/\allowbreak{}readability binding, importgraph cycle check, regression evidence and rollback snapshot. & 22.2 /\allowbreak{} RDV2 & Risk: de-seeding without regression and rollback could make legacy scaffold assumptions appear as current CNNA objects.\\
\hline
\addlinespace[0.24em]
ROADMAP-04 & Modular B core chain and no full-derived QFT face without complement net & Original roadmap chain sheet defines ComplementLocalNet/\allowbreak{}ComplementStateNet/\allowbreak{}ComplementGNS/\allowbreak{}SeparatingProperty/\allowbreak{}TomitaTakesakiData/\allowbreak{}ModularConjugation/\allowbreak{}ComplementKMS/\allowbreak{}ComplementSpectrumCondition/\allowbreak{}BisognanoWichmannSeed and forbids shortcuts from GNS/\allowbreak{}KMS directly to CPT/\allowbreak{}Kramers. & AQFT/\allowbreak{}modular theory is established reference context, not a CNNA input. The plan must derive B objects from A->B handoff and B net/\allowbreak{}state data. & Add B07, Phase 13.8, RDV3/\allowbreak{}22.3 and guardrails GNG47/\allowbreak{}GNG48. & Require ComplementGNS -> SeparatingProperty -> TomitaTakesakiData/\allowbreak{}ModularConjugation -> ComplementKMS -> ComplementSpectrumCondition before BW/\allowbreak{}CPT/\allowbreak{}spin-statistics/\allowbreak{}Kramers claims. & 13.8, 22.3 /\allowbreak{} B07, RDV3 & Risk: bright-only recovery or split-property-like placeholders could be mistaken for a full-derived complement-net QFT face.\\
\hline
\addlinespace[0.24em]
ROADMAP-05 & Nuclearity, DHR/\allowbreak{}DR and charge/\allowbreak{}statistics reconstruction above B core & Original roadmap 6g-6i extension adds NuclearitySeed/\allowbreak{}Nuclearity, DHRSectorsSeed and FieldAlgebraReconstructionSeed plus ChargeStatisticsCoherence. & DHR/\allowbreak{}DR/\allowbreak{}nuclearity are deep established-theory targets;\newline in CNNA they must be comparison/\allowbreak{}consumer gates until source-map-derived profiles exist. & Add B08 and Phase 13.9;\newline keep them after the modular B core and before E matter promotion. & Use nuclearity as phase-space control beyond split property and DHR/\allowbreak{}DR as charge/\allowbreak{}gauge reconstruction bridge, not as imported axioms. & 13.9, 22.3 /\allowbreak{} B08, RDV3 & Risk: too early DHR/\allowbreak{}DR language could smuggle established AQFT reconstruction into CNNA instead of deriving source-compatible gates.\\
\hline
\addlinespace[0.24em]
ROADMAP-06 & ParameterClosure PC-a-PC-d and de-primitive rest parameters & Original roadmap PC-a-PC-d chain requires BranchingWitness, UVSpectralSelector, BranchingSelector, BackreactionSelector, ParameterClosure, HorizonLevel, EffectiveLambda, BackreactionFixedPoint and RegularizationClosure;\newline beta, b and Lmax cannot remain silent primitive exports. & This is a dependency-order/\allowbreak{}gate obligation;\newline no closure theorem is implied by roadmap prose. & Add PC00 and Phase 15.5;\newline add guardrail GNG49. & Route b/\allowbreak{}Lmax/\allowbreak{}beta/\allowbreak{}regularizer handling through derived selectors, horizon/\allowbreak{}fixed-point data or explicit residual-variable markers before AS/\allowbreak{}NCG/\allowbreak{}window language is allowed. & 15.5, 22.4 /\allowbreak{} PC00, RDV4 & Risk: free numerical/\allowbreak{}thermal parameters could remain hidden as legacy inputs while the plan claims derived-only maturity.\\
\hline
\addlinespace[0.24em]
ROADMAP-07 & D/\allowbreak{}E bridge, local covariance, Hadamard/\allowbreak{}Wick and renormalized D language & Original roadmap D minimal tree and D/\allowbreak{}E chain sheet require DerivedSpacetime, InterfaceCausality, LocalCovariance, RelativeCauchyEvolution, DynamicalLocality, HadamardStateSeed, MicrolocalSpectrumGate, LocalWickPolynomialsSeed, DerivedStressTensorSeed and later DHR/\allowbreak{}DR/\allowbreak{}CPT/\allowbreak{}statistics bridges. & LCQFT/\allowbreak{}Hadamard/\allowbreak{}Wick/\allowbreak{}NCG are external target structures;\newline their CNNA role is staged gates/\allowbreak{}comparison layers. & Add D04, DE00, Phases 15.6 and 16.5, and RDV5/\allowbreak{}22.5. & Stage D/\allowbreak{}E so local-covariant and microlocal admissibility precede stress tensor/\allowbreak{}entanglement/\allowbreak{}Einstein/\allowbreak{}NCG comparisons and matter reconstruction. & 15.6, 16.5, 22.5 /\allowbreak{} D04, DE00, RDV5 & Risk: D or E could start from familiar physics vocabulary before the necessary CNNA locality/\allowbreak{}covariance/\allowbreak{}state/\allowbreak{}gate chain exists.\\
\hline
\addlinespace[0.24em]
ROADMAP-08 & Feynman matter audit and semantic cone/\allowbreak{}time-reversal warnings & Original roadmap Feynman audit states: time reversal alone is not matter;\newline anti-linear involution does not replace particle/\allowbreak{}antiparticle dictionaries, spinor representations, CAR/\allowbreak{}Fock structures, exchange signs or chiral inflow;\newline GlobalStateCone is not a lightcone. & Physics interpretation is reference/\allowbreak{}gate material only. Existing time-reversal or cone code cannot be read as full matter/\allowbreak{}causality by name alone. & Add E04, Phase 16.5 and guardrail GNG50. & Require charge/\allowbreak{}anticharge, statistics sign, representation and D/\allowbreak{}E bridge gates before Matter labels are promoted. Keep cone/\allowbreak{}time-reversal semantics disambiguated. & 16.5, 22.6 /\allowbreak{} E04, RDV6 & Risk: word-level overlap such as cone/\allowbreak{}time reversal could create false physical closure in documentation and Explorer views.\\
\hline
\addlinespace[0.24em]
ROADMAP-09 & Late particle/\allowbreak{}scattering/\allowbreak{}IR horizon windows & Original roadmap late horizon layer records ClusterDecompositionSeed, HaagRuelleSeed, ParticleInterpretationWindow and InfraredChargeClassSeed after charge reconstruction and spectral control. & Scattering/\allowbreak{}IR are crucial physics horizons but not immediate core blockers;\newline treat them as windows/\allowbreak{}reference targets. & Add E05 and Phase 16.6. & Keep particle/\allowbreak{}scattering/\allowbreak{}IR visible for long-term completeness while preventing them from blocking Pillar A/\allowbreak{}B/\allowbreak{}C/\allowbreak{}D core closure or entering as primitives. & 16.6, 22.6 /\allowbreak{} E05, RDV6 & Risk: without a horizon row, later phenomenology may either disappear from the plan or be introduced too early as a core object.\\
\hline
\addlinespace[0.24em]
ROADMAP-10 & Meta-audit markers, context sheet and multi-level truth protocol & Original roadmap protocol and meta-audit markers distinguish archive-bound roadmap evidence, formal verification, semantic/\allowbreak{}prose review, physical interpretation, AI as heuristic/\allowbreak{}translation layer, and special windows as windows. & This aligns with current no-infallibility/\allowbreak{}expert-review and Lean/\allowbreak{}YAML/\allowbreak{}tool separation. It still needs explicit ledger/\allowbreak{}checklist integration. & Add Phase 22.7;\newline link RDV0-RDV6 to context-class taxonomy, object significance, total declaration classification and external-source cross-check. & Use meta-markers and context sheets as audit tags, not as mathematical data. Ensure AI/\allowbreak{}prose/\allowbreak{}physics claims are reviewable and do not override kernel-checked code. & 22.7 /\allowbreak{} RDV0-RDV6, CEX11 & Risk: a polished generated PDF/\allowbreak{}Explorer could make roadmap prose appear more certain than checked code or expert-reviewed semantics.\\
\hline
\addlinespace[0.24em]
REALOQS-01 & REALOQS source inventory & Uploaded REALOQS v0.0605 source zip. & Static scan: 270 Lean files, 1727 declarations, 88 PillarB files, 62 PillarB/\allowbreak{}AQFT files, 596 AQFT declarations;\newline no sorry/\allowbreak{}admit/\allowbreak{}axiom hits by text scan. & Accept as legacy code-evidence source only. & Add RLC0 and phase 23.0;\newline archive source and scan reports. & 23.0 /\allowbreak{} RLC0 & No Lean build was run here;\newline text scan does not prove semantic compatibility or current mathlib compatibility.\\
\hline
\addlinespace[0.24em]
REALOQS-02 & AQFT foundation layer & REALOQS PillarB/\allowbreak{}AQFT foundation modules. & StarAlgebra, StateOn, CStar bundles, LocalStarAlgebraNet, StateNet and basic hom/\allowbreak{}restriction lemmas are present. & Useful as target/\allowbreak{}profile map;\newline direct import blocked. & Add 23.1 /\allowbreak{} RLC1 migration-disposition ledger. & 23.1 /\allowbreak{} RLC1 & Some structures may be superseded by mathlib/\allowbreak{}native CNNA abstractions;\newline trivial wrappers and semantic load must be classified.\\
\hline
\addlinespace[0.24em]
REALOQS-03 & A->B dependency inversion & TreeOfCliquesFromPillarA, extended FromA and Meta/\allowbreak{}PillarB files. & Legacy B code often sees old A-side interfaces/\allowbreak{}config/\allowbreak{}harness objects directly. & Treat as incompatibility risk with current CNNA A->B handoff discipline. & Add 23.2 /\allowbreak{} RLC2;\newline require explicit handoff replacement map. & 23.2 /\allowbreak{} RLC2 & Some concepts may still be recoverable, but only after CNNA-derived A->B source data exists.\\
\hline
\addlinespace[0.24em]
REALOQS-04 & BoundaryMatrix AQFT derived layer & REALOQS PillarB/\allowbreak{}AQFT/\allowbreak{}Derived/\allowbreak{}BoundaryMatrix*.lean. & 28 BoundaryMatrix derived modules, 378 declarations, all 28 files marked noncomputable in the scan;\newline many use @[simp]. & Advanced target route;\newline not operational-path compatible as-is. & Add 23.3 /\allowbreak{} RLC3 rewrite and noncomputable-isolation plan. & 23.3 /\allowbreak{} RLC3 & Text scan cannot distinguish harmless noncomputable reference code from semantically essential noncomputability;\newline future module-level review required.\\
\hline
\addlinespace[0.24em]
REALOQS-05 & Haag-Kastler gates & REALOQS PillarB/\allowbreak{}Gates/\allowbreak{}HaagKastler and AQFT/\allowbreak{}HaagKastler modules. & Isotony, locality, additivity, covariance, spectrum, vacuum, time-slice, strong additivity and split-property gate vocabulary exists. & Use as gate checklist and ordering prompt. & Add 23.4 /\allowbreak{} RLC4;\newline route locality/\allowbreak{}spacetime-dependent gates behind C/\allowbreak{}D prerequisites where needed. & 23.4 /\allowbreak{} RLC4 & Legacy predicates may encode assumptions that CNNA must derive later;\newline no closure claim.\\
\hline
\addlinespace[0.24em]
REALOQS-06 & State/\allowbreak{}GNS/\allowbreak{}KMS/\allowbreak{}modular route & REALOQS AQFT State, StateNet, GlobalCone, GNS, KMS, TimeOrientation, TimeReversal and derived KMS/\allowbreak{}GNS/\allowbreak{}Modular modules. & A plausible route state -> state net -> GNS -> dynamics/\allowbreak{}KMS -> modular data is already sketched in code. & Use as sequence target, not imported source. & Add 23.5 /\allowbreak{} RLC5 route comparison. & 23.5 /\allowbreak{} RLC5 & Regime split noOuterEnvironment vs HasOuterEnvironment must be enforced;\newline time reversal alone must not become matter semantics.\\
\hline
\addlinespace[0.24em]
REALOQS-07 & Quasi-local/\allowbreak{}split/\allowbreak{}nuclearity horizon & REALOQS AQFT/\allowbreak{}QuasiLocal and split-property modules plus roadmap B08. & Quasi-local carrier/\allowbreak{}completion/\allowbreak{}star/\allowbreak{}state-lift and split-property themes exist. & Align with later modular B, nuclearity and DHR/\allowbreak{}DR horizon. & Add 23.6 /\allowbreak{} RLC6. & 23.6 /\allowbreak{} RLC6 & May require established AQFT comparison layers;\newline must not be treated as CNNA-generated until B core and source gates exist.\\
\hline
\addlinespace[0.24em]
REALOQS-08 & Examples/\allowbreak{}meta files & REALOQS Examples/\allowbreak{}TreeOfCliques\_AQFT\_* and Meta/\allowbreak{}PillarB\_* files. & Examples and meta closure/\allowbreak{}active-path files are numerous and partly noncomputable. & Quarantine from critical path;\newline mine as regression expectations only. & Add 23.7 /\allowbreak{} RLC7. & 23.7 /\allowbreak{} RLC7 & Examples can overfit old scaffold;\newline regression harnesses need new CNNA source objects before execution.\\
\hline
\addlinespace[0.24em]
ASTRO-27 & Derived-only closure requires generated operational smoke path & Astronews CNNA thread page 2, post \#27. Key content: algebraic idle-run success is insufficient;\newline skeleton, component proof logic and genuine data must all pass through the machine. Lines 237-243 in the 2026-05-12 readout state that only then can one really speak of derived-only. & This is not external theory input;\newline it is a public project self-audit statement by antaris. & Promote to global derived-only closure rule and workflow. & Added POST3, GNG55 and Generated-data derived-only closure workflow. & Global;\newline SM/\allowbreak{}AR/\allowbreak{}AQFT/\allowbreak{}Explorer status logic;\newline CEX0/\allowbreak{}CEX11. & Forum text is interpretive project governance, not Lean evidence. It constrains status semantics but does not prove any object.\\
\hline
\addlinespace[0.24em]
ASTRO-27B & Missing structures must be generated at their natural source layer & Astronews CNNA thread page 2, post \#27 lines 240-243: repeated downstream failures require jumping back through earlier phases until the natural emergence site is found;\newline forcing the missing structure at the failure site created chaos in the first attempt. & Consistent with existing CNNA provenance rule but stronger: the origin-selection decision itself must be documented. & Promote to red-blocker/\allowbreak{}natural-origin workflow and no-go rule. & Added GNG56 and Natural-origin missing-structure workflow;\newline links to pre-implementation red status and addressing-phase generation. & All implementation phases;\newline pre-implementation phase-check policy;\newline lifecycle/\allowbreak{}provenance ledgers. & The natural origin may require expert review and may not be decidable from a text scan alone.\\
\hline
\addlinespace[0.24em]
ASTRO-28 & YAML/\allowbreak{}Explorer as epistemic control layer beyond linear PDF/\allowbreak{}TeX manageability & Astronews CNNA thread page 2, post \#28 lines 288-307: generalized reusable smoke-test/\allowbreak{}generator, implementation thought-experiment checks, PDF/\allowbreak{}TeX insufficiency, YAML as machine-readable living object, SQL/\allowbreak{}MySQL website and Mindmap/\allowbreak{}Universe-Explorer vision. & This is a project-tooling self-audit and UI/\allowbreak{}knowledge-management requirement. & Promote to Explorer motivation, thought-experiment readiness workflow and generated-secondary documentation rule. & Added GNG57, GNG58, Implementation-thought-experiment phase-readiness workflow and Explorer/\allowbreak{}YAML epistemic-control workflow. & 21.0-21.9;\newline CEX0-CEX11;\newline PIPC policy;\newline SQL/\allowbreak{}Explorer future mirror. & The UI/\allowbreak{}database realization is future infrastructure;\newline YAML remains planning source and Lean remains truth source.\\
\hline
\addlinespace[0.24em]
ASTRO-29 & Foundation no-free-parameter axiomatics and regime split & Astronews CNNA thread page 2, post \#29 lines 330-341: Foundation has only incoming dependency arrows, exposes the no-free-parameter foundation;\newline axioms are root/\allowbreak{}empty-set infinite partial order, finite subsystems, root approximation closed/\allowbreak{}unitary/\allowbreak{}no dark-interface, non-root approximation complement-bearing/\allowbreak{}open bright/\allowbreak{}interface/\allowbreak{}dark split;\newline L\_max/\allowbreak{}horizon and branching require derivation/\allowbreak{}backreaction. & This sharpens the global foundation/\allowbreak{}regime semantics already used in Pillar A/\allowbreak{}B planning. & Promote to global postulate and parameter provenance guardrails. & Added POST2, GNG59, GNG60 and Pillar co-development/\allowbreak{}acyclic handoff workflow. & Foundation;\newline ParameterClosure;\newline A->B/\allowbreak{}A->C/\allowbreak{}A->D/\allowbreak{}A->E handoffs;\newline Pillar B-E planning. & Forum statement is governance/\allowbreak{}interpretation;\newline Lean modules still need to prove each formal claim at the appropriate layer.\\
\hline
\addlinespace[0.24em]
ASTRO-29B & Pillar B-E co-development without provenance cycles & Astronews CNNA thread page 2, post \#29 lines 338-341: L\_max likely needs Pillar D spacetime reconstruction;\newline D may need A, B, C and parts of E;\newline B-E may need deep simultaneous derivation from A;\newline no current need for stochastic quantities. & This challenges a naive linear A->B->C->D->E order but does not authorize cyclic proof dependencies. & Represent co-development with explicit acyclic handoff/\allowbreak{}backreaction windows and edge semantics. & Added GNG61 and Pillar co-development and acyclic handoff workflow. & B-E sidechains;\newline CNNA Explorer graph edge semantics;\newline handoff ledgers. & Exact handoff topology remains open until future phases derive the relevant objects.\\
\hline
\addlinespace[0.24em]
\end{longtable}
\arrayrulecolor{black}
\normalsize
